\documentclass[aspectratio=169,11pt]{beamer}

\usetheme{Madrid}
\usecolortheme{default}
\setbeamertemplate{navigation symbols}{}

\usepackage[T1]{fontenc}
\usepackage[utf8]{inputenc}
\usepackage{lmodern}
\usepackage{microtype}
\usepackage{graphicx}
\usepackage{booktabs}
\usepackage{tabularx}
\usepackage{array}
\usepackage{hyperref}

\definecolor{PGNavy}{HTML}{1E2A44}
\definecolor{PGTeal}{HTML}{2D6F73}
\definecolor{PGGold}{HTML}{C78A2C}
\definecolor{PGLight}{HTML}{F5F1E8}
\definecolor{PGRed}{HTML}{8F2D2D}

\setbeamercolor{palette primary}{bg=PGNavy, fg=white}
\setbeamercolor{palette secondary}{bg=PGTeal, fg=white}
\setbeamercolor{palette tertiary}{bg=PGGold, fg=white}
\setbeamercolor{structure}{fg=PGNavy}
\setbeamercolor{frametitle}{bg=PGLight, fg=PGNavy}
\setbeamercolor{block title}{bg=PGNavy, fg=white}
\setbeamercolor{block body}{bg=PGLight, fg=black}
\setbeamercolor{alerted text}{fg=PGRed}

\title{Functional Assessment, Punishment, and Restorative Possibility}
\subtitle{EDU 442 - Week 8 discussion presentation}
\author{Piter Garcia}
\date{March 2026}

\begin{document}

\begin{frame}
  \titlepage
\end{frame}

\begin{frame}{Working claim}
\begin{block}{Main argument}
Functional assessment is never neutral. If school discipline is already shaped by punishment, surveillance, and narrow ideas of the ``good student,'' then behavior assessment can easily become another tool for control instead of a tool for understanding.
\end{block}

\vspace{0.4cm}
\begin{itemize}
  \item Week 8 pushes me to ask what behavior means inside a racialized punishment system.
  \item The readings move from history (\textbf{Stevenson}) to school structure (\textbf{James-Gallaway}) to implementation (\textbf{Marsh}).
  \item My goal is not to reject assessment, but to reframe it toward context, access, and repair.
\end{itemize}
\end{frame}

\begin{frame}{Why this week matters}
\begin{columns}[T]
  \column{0.52\textwidth}
  \begin{block}{What the week is about}
  \begin{itemize}
    \item School-to-prison pipeline
    \item Discipline as part of a larger punishment regime
    \item Restorative practice as culture change, not a quick fix
  \end{itemize}
  \end{block}

  \column{0.45\textwidth}
  \begin{block}{Question I am carrying}
  When schools label behavior as ``defiance,'' whose norms are being protected, whose needs are being ignored, and what gets misread as a student problem?
  \end{block}
\end{columns}
\end{frame}

\begin{frame}{Stevenson (2021): punishment is structural}
\begin{itemize}
  \item Stevenson situates punishment in the afterlife of slavery and the 13th Amendment's punishment clause.
  \item The point is not just that punishment exists, but that it has been racialized from the start.
  \item In schools, this matters because discipline systems often look administrative or neutral while carrying carceral logics.
  \item If we only count referrals, suspensions, or ``incidents,'' we miss the history producing those categories.
\end{itemize}

\vspace{0.3cm}
\begin{alertblock}{My takeaway}
Discipline data is not raw truth. It is produced inside a system that already decides whose behavior looks threatening, disruptive, or punishable.
\end{alertblock}
\end{frame}

\begin{frame}{James-Gallaway (2022): the ``good student'' as compliance}
\begin{itemize}
  \item James-Gallaway shows how No Excuses discipline normalizes surveillance, control, and dehumanization.
  \item ``Good student'' gets defined as quiet, compliant, and easily managed.
  \item That standard is racialized and ableist: it punishes movement, affect, refusal, self-protection, and different ways of being.
  \item Resistance then gets read as proof that harsher control is necessary.
\end{itemize}

\vspace{0.3cm}
\begin{block}{Why this matters for assessment}
If behavior is assessed against a compliance norm, then the assessment is already biased before any form is filled out.
\end{block}
\end{frame}

\begin{frame}{Marsh (2017): restorative practice needs culture change}
\begin{itemize}
  \item Marsh argues restorative practice fails when schools treat it like a program layered on top of punishment.
  \item Real restorative work requires time, trust, training, and changes in adult behavior.
  \item It is not just a circle, a script, or a reduction in office referrals.
  \item Success should include stronger relationships, more student voice, better repair, and less harm.
\end{itemize}

\vspace{0.3cm}
\begin{alertblock}{Implementation point}
If the institution stays punitive, restorative language can become another form of control.
\end{alertblock}
\end{frame}

\begin{frame}{My lens: why this feels personal}
\begin{columns}[T]
  \column{0.43\textwidth}
  \centering
  \includegraphics[width=\linewidth]{../EDU442_Self_Portrait_Garcia/self_portrait-image0.png}

  \column{0.53\textwidth}
  \begin{itemize}
    \item In my self-portrait, I name myself as a neurodivergent Dominican-American educator and researcher.
    \item I describe institutions that pathologize culture, language, disability, and gender variance rather than understand them.
    \item That is why week 8 lands hard for me: behaviors tied to ADHD, chronic illness, trauma, or cultural difference are often read as noncompliance.
    \item I am interested in equitable diagnostics because I have lived diagnostic injustice.
  \end{itemize}
\end{columns}
\end{frame}

\begin{frame}{Reframing functional assessment}
\small
\begin{tabularx}{\textwidth}{>{\raggedright\arraybackslash}p{0.30\textwidth} >{\raggedright\arraybackslash}X}
\toprule
\textbf{Punitive reading} & \textbf{Justice-oriented functional assessment} \\
\midrule
``Off task'' & What in the task, environment, pacing, or sensory load made engagement hard? \\
``Defiant'' & Was the student protecting autonomy, dignity, regulation, or safety? \\
``Disrespectful'' & Are we misreading cultural language, tone, emotion, or disability-related communication? \\
``Noncompliant'' & Were expectations explicit, accessible, and reasonable in the first place? \\
``Behavior problem'' & What adult routines, school norms, or structural pressures are escalating this pattern? \\
\bottomrule
\end{tabularx}

\vspace{0.3cm}
\normalsize
\begin{block}{My claim}
Functional assessment should ask what the behavior does in context, but also who built the context and who benefits from its norms.
\end{block}
\end{frame}

\begin{frame}{What accountable practice would look like}
\begin{itemize}
  \item Start with an environment audit before blaming the student: task clarity, sensory load, transition design, relational safety.
  \item Include student voice and family knowledge, not just adult interpretation.
  \item Distinguish harm from norm violation: not every rule-breaking act is dangerous.
  \item Measure more than fewer referrals: belonging, participation, repair, trust, and re-entry after conflict.
  \item Resource the work: restorative approaches cannot run on under-training and goodwill alone.
\end{itemize}
\end{frame}

\begin{frame}{My closing synthesis}
\begin{block}{What I am taking from Week 8}
Stevenson helps me see punishment as historical and structural. James-Gallaway shows how that structure enters school through surveillance and compliance. Marsh reminds me that alternatives only matter if adults and institutions actually change.
\end{block}

\vspace{0.4cm}
\begin{alertblock}{Bottom line}
If assessment is disconnected from power, it will misname domination as data. If assessment is tied to context, dignity, and repair, it can become part of justice-oriented teaching.
\end{alertblock}
\end{frame}

\begin{frame}{Discussion questions}
\begin{enumerate}
  \item Where do school rules construct the ``good student'' as compliance, and how does that punish neurodivergence, trauma responses, or disability when behaviors are read as defiance?
  \item What would accountable restorative implementation look like in an under-resourced school under heavy surveillance, and what should we measure besides fewer referrals?
\end{enumerate}
\end{frame}

\begin{frame}{References}
\footnotesize
\begin{itemize}
  \item James-Gallaway, C. D. (2022). ``The kids in prison program'': A critical race personal counternarrative.
  \item Marsh, V. (2017). Becoming restorative: Three schools transitioning to a restorative practices culture.
  \item Stevenson, B. (2021). ``Punishment.'' In \textit{The 1619 Project}, pp. 431--456.
  \item DuVernay, A. (Director). (2016). \textit{13th}.
  \item Garcia, P. (2026). \textit{The Puzzle Plan: A Self-Portrait in Intersectionality}.
\end{itemize}
\end{frame}

\end{document}
